\documentclass[a4paper,11pt]{article}

% --------------------------------------------------
% Packages
% --------------------------------------------------
\usepackage[utf8]{inputenc}
\usepackage[T1]{fontenc}
\usepackage[english]{babel}

\usepackage[
    backend=biber,
    style=numeric,
    natbib=true
]{biblatex}
\addbibresource{references.bib}

\usepackage{geometry}
\geometry{margin=2.5cm}

\usepackage{graphicx}
\graphicspath{{./images/}}

\usepackage{caption}

\usepackage{booktabs}
\usepackage{amsmath}
\usepackage{amssymb}
\usepackage{float}
\usepackage{subcaption}
\usepackage{siunitx}
\usepackage{xcolor}
\usepackage{enumitem}

\usepackage{hyperref}
\hypersetup{
    colorlinks=true,
    linkcolor=blue,
    citecolor=blue,
    urlcolor=blue
}

\usepackage{titlesec}

\renewcommand{\thesection}{\Alph{section}}
\titleformat{\section}
  {\normalfont\Large\bfseries}
  {Part \thesection\ ---}
  {0.5em}
  {}

\usepackage{float}
\usepackage{placeins}

\widowpenalty=10000
\clubpenalty=10000

% --------------------------------------------------
% Title
% --------------------------------------------------
\title{
    \textbf{Reinforcement Learning}\\
    \large Project: Cloud Auto-Scaling with Deep Reinforcement Learning
}

\author{
    Loïc Thin \\
    Jonas Keller\\
}

\date{\today}

% --------------------------------------------------
% Document
% --------------------------------------------------
\begin{document}

\maketitle

\tableofcontents
\newpage


% ==================================================
\section{Understand the system and establish baselines}
% ==================================================
\subsection{Dynamics of server startup and scale-down}
The given environment consists of two parts: the server management and the request processing.

\subsection{MDP formalization}
We model the problem with the MDP $\mathcal{M} = \left(\mathcal{S}, \mathcal{A}, P, R, \gamma \right)$.
At time step $t$, the state $s_t \in \mathcal{S}$ is defined by
\begin{equation*}
    s_t = \left(
        n_t, q_t, d_t, d_{t-1},
        \sin(\theta_t), \cos(\theta_t),
        p_{1,t}, \ldots, p_{K,t}
    \right)
\end{equation*}
where the individual components are as follows:
\begin{table}[H]
    \centering
    \begin{tabular}{ll}
        \hline
        \textbf{Component} & \textbf{Description} \\
        \hline
        $n_t$ & Normalized number of active servers. \\
        $q_t$ & Normalized queue length. \\
        $d_t$ & Current normalized incoming demand. \\
        $d_{t-1}$ & Previous normalized incoming demand. \\
        $\sin(\theta_t), \cos(\theta_t)$ & Periodic encoding of the time of day. \\
        $p_{k,t}$ & Normalized number of pending servers with $k$ steps remaining. \\
        \hline
    \end{tabular}
    \caption{Components of the state representation.}
    \label{tab:state_components}
\end{table}

\noindent
The action space is given by $\mathcal{A} = \{0,1,2\}$ with the following semantics
\begin{table}[H]
    \centering
    \begin{tabular}{ll}
        \hline
        \textbf{Action code} & \textbf{Description} \\
        \hline
        0 & Remove one active server \\
        1 & Hold \\
        2 & Start one server \\
        \hline
    \end{tabular}
    \caption{Semantics of the action codes.}
    \label{tab:action_codes}
\end{table}

Next, we can formalize the transition mechanisms. Starting from a state $s_t \in \mathcal{S}$, by performing an action $a_t \in \mathcal{A}$, we can describe the next state $s_{t+1}$ or rather its conditional probability distribution $P(s_{t+1}|s_t,a_t)$.
The transition itself consists of three mechanisms that contribute to the resulting next state.

Firstly, there are the action-dependent updates. The number of active servers and pending servers changes according to the selected action and the deterministic startup process.
It starts by decreasing the remaining steps for each pending server.
This means that we get
\begin{equation*}
    p_{k-1,t+1} = p_{k,t}\text{,}\quad k\in {2, ... K} \\
\end{equation*}
which simply shifts the normalized number of pending servers by one to the next lower count of remaining steps.

The pending servers $p_{1,t}$, which previously only needed one more time step, are now added to the active servers from the previous step, resulting in
\begin{equation*}
    n_{t+1} = n_t + p_{1,t} - \begin{cases}
        1 / N_{max}\text{,}\quad \text{if} \quad a_t=0 \\
        0\text{,}\quad \text{otherwise} 
    \end{cases}
\end{equation*}
where $N_{max}$ is a hyperparameter denoting the maximum number of servers available to our application.

In addition to that, we subtract one active server if $a_t=0$  or add one pending server with the full $K$ steps remaining to the pending server list if $a_t = 2$ (and both are valid actions).
This yields
\begin{equation*}
    p_{K,t+1} = \begin{cases}
        1\text{,}\quad \text{if} \quad a_t=2 \\
        0\text{,}\quad \text{otherwise} 
    \end{cases}
\end{equation*}

Secondly, fixed, rule-based system dynamics evolve the queue according to the incoming demand and the current server capacity.
Therefore, $q_{t+1}$ only depends on $n_{t+1}$, $q_t$ and $d_{t+1}$ and is updated as follows
\begin{equation*}
    q_{t+1} = 
\end{equation*}

% ==================================================
\subsection{Evaluation of fixed-capacity and reactive-threshold policies}
% ==================================================

We evaluate the two supplied non-learning policies using three workload
profiles: \texttt{daily}, \texttt{volatile} and \texttt{shifted\_peak}.
For each profile, the policies are evaluated over 10 episodes using common
random seeds. This paired evaluation makes the comparison more reliable,
since each policy is exposed to the same set of episode seeds.

The first policy, \texttt{fixed5}, maintains a fixed capacity of five active
servers. It provides a simple reference point, but it cannot adapt its
capacity to changes in demand. The second policy, \texttt{threshold}, is the
provided reactive-threshold controller. Unlike the fixed policy, it adjusts
the number of servers according to the observed system conditions.

We also implement a stronger non-learning baseline, named
\texttt{personal}. It uses a three-observation moving average of the total
observed demand, defined as incoming arrivals plus the current queue length.
It applies two distinct thresholds to avoid reacting identically to demand
increases and decreases:
\begin{itemize}
    \item The scale-up threshold is set to $0.8$. A new server is requested
    when the moving average exceeds $0.8$ times the available capacity.
    Available capacity accounts for both active and pending servers.
    Consequently, the controller does not request additional servers when
    enough capacity is already starting up.
    \item The scale-down threshold is set to $0.35$. An active server is
    removed when the moving average falls below $0.35$ times the capacity
    of the active servers, provided that the minimum server constraint is
    respected.
\end{itemize}
If neither condition is met, the controller keeps the current configuration.
The use of separate upper and lower thresholds introduces hysteresis, while
the moving average reduces sensitivity to short-lived demand fluctuations.
Taking pending servers into account when scaling up also avoids redundant
server-start requests. This policy therefore combines a simple traffic-trend
smoother, hysteresis and pending-server awareness, without using learning.

\subsection{Reward and operational metrics}
% ==================================================

A higher reward alone does not necessarily indicate a better autoscaling
policy. The reward is an aggregate objective, whereas cost, drop rate and
the number of configuration switches describe distinct operational aspects
of the system. We therefore report the reward together with these raw
metrics. All values below are means and standard deviations over 10 episodes.

\begin{table}[H]
    \centering
    \small
    \begin{tabular}{lrrrr}
        \toprule
        \textbf{Policy} &
        \textbf{Reward} &
        \textbf{Cost} &
        \textbf{Drop rate} &
        \textbf{Switches} \\
        \midrule
        fixed5 &
        $-109.49 \ (4.49)$ &
        $86.26 \ (0.00)$ &
        $0.01276 \ (0.00279)$ &
        $2.00 \ (0.00)$ \\
        threshold &
        $-67.72 \ (2.20)$ &
        $67.68 \ (2.22)$ &
        $0.00000 \ (0.00000)$ &
        $29.70 \ (5.00)$ \\
        personal &
        $\mathbf{-63.81} \ (2.27)$ &
        $\mathbf{63.32} \ (2.47)$ &
        $0.00000 \ (0.00000)$ &
        $\mathbf{17.40} \ (2.06)$ \\
        \bottomrule
    \end{tabular}
    \caption{Evaluation results on the daily workload profile.}
    \label{tab:daily_baselines}
\end{table}

\begin{table}[H]
    \centering
    \small
    \begin{tabular}{lrrrr}
        \toprule
        \textbf{Policy} &
        \textbf{Reward} &
        \textbf{Cost} &
        \textbf{Drop rate} &
        \textbf{Switches} \\
        \midrule
        fixed5 &
        $-107.52 \ (10.60)$ &
        $86.26 \ (0.00)$ &
        $0.01170 \ (0.00660)$ &
        $2.00 \ (0.00)$ \\
        threshold &
        $-71.20 \ (3.13)$ &
        $70.23 \ (2.83)$ &
        $0.00000 \ (0.00000)$ &
        $82.70 \ (11.93)$ \\
        personal &
        $\mathbf{-67.71} \ (3.95)$ &
        $\mathbf{64.50} \ (2.70)$ &
        $0.00029 \ (0.00044)$ &
        $\mathbf{47.30} \ (9.60)$ \\
        \bottomrule
    \end{tabular}
    \caption{Evaluation results on the volatile workload profile.}
    \label{tab:volatile_baselines}
\end{table}

\begin{table}[H]
    \centering
    \small
    \begin{tabular}{lrrrr}
        \toprule
        \textbf{Policy} &
        \textbf{Reward} &
        \textbf{Cost} &
        \textbf{Drop rate} &
        \textbf{Switches} \\
        \midrule
        fixed5 &
        $-111.53 \ (4.43)$ &
        $86.26 \ (0.00)$ &
        $0.01404 \ (0.00290)$ &
        $2.00 \ (0.00)$ \\
        threshold &
        $-67.98 \ (1.64)$ &
        $67.95 \ (1.63)$ &
        $0.00000 \ (0.00000)$ &
        $31.10 \ (4.59)$ \\
        personal &
        $\mathbf{-64.32} \ (2.09)$ &
        $\mathbf{63.56} \ (2.21)$ &
        $0.00000 \ (0.00000)$ &
        $\mathbf{19.90} \ (1.92)$ \\
        \bottomrule
    \end{tabular}
    \caption{Evaluation results on the shifted-peak workload profile.}
    \label{tab:shifted_peak_baselines}
\end{table}

The results show that \texttt{fixed5} performs substantially worse in terms
of reward and cost than the adaptive policies. Its fixed capacity cannot
adjust to the workload, and it also exhibits non-zero request drop rates
on all three profiles. Although it performs only two configuration switches
per episode, this low number should not be interpreted as better control:
it is a consequence of keeping a fixed capacity rather than adapting to
demand.

The \texttt{personal} controller obtains a better mean reward than
\texttt{threshold} on all three workload profiles. Since rewards are
negative, a value closer to zero is better. The improvement is accompanied
by a lower mean cost: approximately $6.4\%$ on \texttt{daily},
$8.2\%$ on \texttt{volatile} and $6.5\%$ on \texttt{shifted\_peak},
relative to \texttt{threshold}. Thus, the higher reward is not obtained
by spending more resources.

The number of configuration switches is also lower for
\texttt{personal} in every profile: it decreases from $29.7$ to $17.4$
on \texttt{daily}, from $82.7$ to $47.3$ on \texttt{volatile}, and from
$31.1$ to $19.9$ on \texttt{shifted\_peak}. This is consistent with the
intended effect of the moving average and hysteresis, which reduce
unnecessary reactions to short-term variations in demand.

Regarding service quality, both adaptive policies have a zero mean drop
rate on the \texttt{daily} and \texttt{shifted\_peak} profiles. On the
\texttt{volatile} profile, \texttt{personal} has a small non-zero mean
drop rate of $0.00029$ (approximately $0.029\%$), while
\texttt{threshold} has no observed drops. This is a slight degradation
that should not be hidden by the reward improvement. Its practical
acceptability depends on the service-level constraint imposed on request
losses; the evaluation results alone do not define such an acceptable
threshold. The reported standard deviation of the drop rate for
\texttt{personal} on this profile is $0.00044$, which also indicates
variability between episodes.

Overall, the custom \texttt{personal} baseline provides the best balance
among the tested non-learning policies: it consistently improves reward
and reduces cost and configuration switches compared with the supplied
reactive-threshold policy. The only observed operational disadvantage is
the very small drop rate on the volatile workload. This result should
therefore be considered alongside the application's acceptable request-loss
limit, rather than relying on reward alone.

\end{document}
